\begin{frame}{자주 묻는 질문 (2/2)}
  \textbf{Q. 서포트 벡터의 개수가 $C$에 따라 변하나?}
  \begin{itemize}
    \item C 스윕 표에서: $C=0.01$ $\to$ 14개, $C=0.1$ $\to$ 5개,
      $C\ge1$ $\to$ 2개.
    \item \textbf{$C$가 작을수록 SV가 많아짐} — 이상치를 무시하고
      마진을 넓히려면 더 많은 점이 마진 경계에 달라붙어야 하므로.
    \item SV 개수 = SVM의 \textbf{효과적 모델 크기}(예측 시
      비교할 점의 수) — $C$ 튜닝의 부수적 효과로 ``모델의 간결함''도
      함께 조절.
  \end{itemize}
  \vskip0.5em
  \textbf{Q. 이상치가 \textbf{많이} 있으면 $C$ 튜닝으로 해결되지 않나요?}
  \begin{itemize}
    \item 이상치가 \textbf{데이터의 상당 부분}(예: 30\% 이상)이면
      어느 쪽을 택해도 일반화 성능이 떨어짐 — 작은 $C$는 진짜
      신호도 함께 무시, 큰 $C$는 노이즈에 과적합.
    \item 이런 때는 \textbf{모델이 아니라 데이터 자체}의 문제를 먼저:
      라벨링 품질 점검, 이상치 탐지(4장 kNN, 8장 앙상블),
      \textbf{주변부 잘라내기}(클러스터 밖에 고립된 의심점 전처리
      제거).
    \item $C$는 ``몇 \% 수준의 노이즈''를 조절하는 손잡이지,
      ``데이터의 절반이 틀린'' 문제를 해결하는 \textbf{마법이
      아니다}.
  \end{itemize}
\end{frame}
